\section{Werkstoffstufen, Dehnraten und Radkörper}
\label{sec:stufen_raten_koerper}

Dieser Abschnitt trennt am Sektormodell aus \secref{sec:fe_modell} die Einflüsse der Werkstoffstufe, der Dehnrate der Werkstoffkarte und des Radkörpers. Grundlage ist die Serie des Ringkörpers mit 99 Stellungen von einer Teilung vor A bis einer Teilung hinter E mit der quasistatischen Karte (QS). Gerechnet wurden die vier Stufen W1 (linear-elastisch, isotrop), W2 (elastisch-plastisch, isotrop), W3 (linear-elastisch, anisotrop) und W4 (elastisch-plastisch, anisotrop) nach der Abstufung von \cite{Sto18}; 395 der 396 Rechnungen liefen vollständig durch (vgl. \secref{sec:kopfkante}). Die Raten DY1 (Dehnrate $3$\,s$^{-1}$, $E = 2543$\,MPa) und DY2 ($60$\,s$^{-1}$, $E = 3597$\,MPa) wurden an den Normpunkten A bis E gerechnet, ebenso der Taschenkörper nach Zeichnung. \Note{PH: Elastizitätsmodul der QS-Karte bei 80 Grad C zum Vergleich (Werkstoffkarte, Anhang D).} Alle Werte sind Größtwerte über die Zahnbreite; $\sigma_{\mathrm{1}}$ bezeichnet die größte Zugspannung und $|\sigma_{\mathrm{3}}|$ die größte Druckspannung in den Fußrundungen, der Kontaktdruck den größten Wert auf der Radflanke.

Die Anisotropie der linearen Stufe hebt die Fußspannungen an und lässt die Verformung nahezu unverändert (\tabref{tab:stufen_16nm}). Bei $16$\,Nm liegt der Größtwert von $\sigma_{\mathrm{1}}$ in W3 um +$17{,}7$\,\% über W1, an den Punkten A bis E um +$15{,}8$\,\% bis +$19{,}6$\,\%; $|\sigma_{\mathrm{3}}|$ nimmt um +$13{,}9$\,\% zu. Die Ritzeldrehung wächst an den Normpunkten nur um +$1{,}8$\,\% bis +$4{,}6$\,\%, die Kopfverschiebung ändert sich um höchstens $2{,}6$\,\%. Der Kontaktdruck reagiert je nach Stellung verschieden; an C sinkt er um $11{,}5$\,\%, an E steigt er um $14{,}5$\,\%.

Die Plastizität wirkt in die Gegenrichtung und deutlich stärker. In W2 liegt der Größtwert von $\sigma_{\mathrm{1}}$ bei $16$\,Nm mit $52{,}89$\,MPa um $38{,}5$\,\% unter dem Wert von W1 mit $86{,}03$\,MPa, $|\sigma_{\mathrm{3}}|$ sinkt um $42{,}6$\,\% und der größte Kontaktdruck um $46{,}0$\,\%. Zugleich nimmt die Ritzeldrehung um +$17{,}4$\,\% zu, an den Punkten A, B, D und E um +$22{,}6$\,\% bis +$26{,}6$\,\%, an C um +$12{,}9$\,\%. Schon bei $8$\,Nm beträgt der Abstand in $\sigma_{\mathrm{1}}$ $31{,}9$\,\%, weil die Fließkurve der Karte bei $80$\,\textdegree C bereits bei $4{,}3$\,MPa beginnt. Zwischen W3 und W4 wiederholt sich das Bild mit $-$$33{,}2$\,\% in $\sigma_{\mathrm{1}}$, $-$$35{,}0$\,\% in $|\sigma_{\mathrm{3}}|$ und +$13{,}4$\,\% in der Drehung; der Kontaktdruck an C steigt in W4 jedoch um $18{,}1$\,\%. Anisotropie und Plastizität überlagern sich nahezu multiplikativ, wie der Vergleich von W4 mit W2 zeigt. Die Fußspannungen liegen dort um +$27{,}9$\,\% ($\sigma_{\mathrm{1}}$) und +$29{,}1$\,\% ($|\sigma_{\mathrm{3}}|$) höher, die Drehung ist mit $-$$0{,}7$\,\% bis +$1{,}7$\,\% praktisch gleich; die Faserorientierung bleibt also auch in der plastischen Stufe maßgebend.

\begin{table}[htbp]
\centering
\caption{Werkstoffstufen am Ringkörper bei $16$\,Nm, quasistatisch; Größtwerte der Stufe W1 über 99 Stellungen und relative Änderung der Größtwerte zwischen den Stufen.}
\label{tab:stufen_16nm}
\footnotesize
\setlength{\tabcolsep}{4pt}
\begin{tabular}{L{4.4cm}lllll}
\hline
Vergleich & Drehung & Kopfweg & Kontaktdruck & $\sigma_{\mathrm{1}}$ & $|\sigma_{\mathrm{3}}|$ \\
\hline
W1, Größtwert & $3{,}54$\,mrad & $101{,}40$\,\textmu m & $184{,}58$\,MPa & $86{,}03$\,MPa & $107{,}50$\,MPa \\
W3 gegen W1 (Anisotropie, linear) & +$2{,}8$\,\% & $-$$2{,}3$\,\% & +$10{,}3$\,\% & +$17{,}7$\,\% & +$13{,}9$\,\% \\
W2 gegen W1 (Plastizität, isotrop) & +$17{,}4$\,\% & +$13{,}8$\,\% & $-$$46{,}0$\,\% & $-$$38{,}5$\,\% & $-$$42{,}6$\,\% \\
W4 gegen W3 (Plastizität, anisotrop) & +$13{,}4$\,\% & +$13{,}4$\,\% & $-$$46{,}1$\,\% & $-$$33{,}2$\,\% & $-$$35{,}0$\,\% \\
W4 gegen W2 (Anisotropie, plastisch) & $-$$0{,}7$\,\% & $-$$2{,}7$\,\% & +$10{,}1$\,\% & +$27{,}9$\,\% & +$29{,}1$\,\% \\
\hline
\end{tabular}
\end{table}
% Bildwunsch: 80_output/fe/kst_e/series_20layers_bore33_round0p1f12_30perpitch_margin1_qs_gerechnet/analyse/vergleich_stufen_LOAD_16NM.svg, Stufenvergleich W1 bis W4 über dem Wälzweg (Ringkörper, 16 Nm, QS); überschneidet sich mit dem Bildplan von Abschnitt 5.2, dort oder hier.

Die Dehnrate wirkt vor allem auf die Verformung. Mit DY1 sinkt die Ritzeldrehung an A bis E bei $16$\,Nm je nach Stufe um $5{,}3$\,\% bis $10{,}1$\,\%, $\sigma_{\mathrm{1}}$ steigt um $0{,}1$\,\% bis $9{,}0$\,\%, am stärksten in W4. Mit DY2 geht die Drehung an A, B, D und E um $34{,}0$\,\% bis $41{,}0$\,\% zurück, an C um $22{,}6$\,\% bis $26{,}2$\,\%. An A, B, D und E ändert sich $\sigma_{\mathrm{1}}$ unter DY2 bei den linearen Stufen um höchstens $1{,}9$\,\%, bei W2 und W4 um +$7{,}8$\,\% bis +$12{,}1$\,\%. An C dagegen steigt $\sigma_{\mathrm{1}}$ in allen vier Stufen um +$21{,}2$\,\% bis +$27{,}4$\,\%. Der Kontaktdruck an C nimmt in den linearen Stufen um +$35{,}9$\,\% und +$37{,}2$\,\% zu, über alle Lastschritte um bis zu $40{,}2$\,\%. W2 und W4 zeigen dort bei $16$\,Nm +$3{,}0$\,\% und $-$$13{,}9$\,\%. Eine Erklärung bietet der verlängerte Eingriff der Serie. Der mittlere Zahn trägt schon $0{,}47$ bis $0{,}60$ Teilungen vor A und noch $0{,}35$ bis $0{,}48$ Teilungen hinter E. In der Stellung C übernehmen deshalb auch die Nachbarzähne Last. Ihnen teilt der steifere Werkstoff weniger zu, und der mittlere Zahn trägt mehr, obwohl sich das Paar weniger verdreht.
% Bildwunsch: 80_output/fe/kst_e/pilot_20layers_bore33_round0p1f12_dy1dy2_AtoE_gerechnet/analyse/vergleich_referenz_LOAD_16NM.svg, Ratenvergleich DY1 und DY2 gegen QS an A bis E (Ringkörper, 16 Nm).
% Hinweis an den Nutzer: Tagebuch (outcome des Ratenpiloten) und ADR-116 Punkt 47(f) nennen den DY2-Anstieg von sigma1 um 20 bis 32 % "an A"; die Tabellen in Abschnitt 8 der Analyse zeigen ihn an C (Spaltenreihenfolge Maximum, A, B, C, D, E; an den Absolutwerten des Körpervergleichs bei 8 Nm geprüft). Der Text folgt den Tabellen.

Der Taschenkörper nach Zeichnung unterscheidet sich vom Ringkörper in der Geometrie und in der Schichtung über der Zahnbreite. Beim Ring liegen 20 Schichten zu $0{,}75$\,mm übereinander. Beim Taschenkörper sind es sechs Schichten zu $0{,}667$\,mm über dem Steg und sieben Schichten zu $0{,}786$\,mm je Taschentiefe, so dass CONVERSE die Faserorientierung an anderen Ebenen abtastet. Um beide Anteile zu trennen, wurde der Ringkörper zusätzlich auf den Schichten des Taschenkörpers gerechnet. Der Schichteffekt vergleicht diesen Ring mit dem Ring der Serie, gleiche Geometrie bei anderer Abtastung der Orientierung. Dagegen vergleicht der Körpereffekt den Taschenkörper mit dem Ring gleicher Schichtung, gleiche Abtastung bei anderer Geometrie. Beide Körper liegen für W1 bis W4 und die drei Raten an A bis E vor, je 60 Rechnungen, alle vollständig.

Der Schichteffekt bleibt bei den isotropen Stufen unter $0{,}2$\,\% (\tabref{tab:koerper_spannen}). Bei W3 und W4 erreicht er im Kontaktdruck $-$$5{,}8$\,\% bis +$7{,}1$\,\% und in $|\sigma_{\mathrm{3}}|$ $-$$1{,}3$\,\% bis +$3{,}6$\,\%, in Drehung und $\sigma_{\mathrm{1}}$ höchstens $0{,}8$\,\%. Die andere Abtastung dreht die lokale Faserrichtung bei $21$\,\% der Zahnelemente um mehr als $10$\textdegree; das erreicht den Kontakt der anisotropen Stufen, nicht aber die Verformung. Jeder anisotrope Vergleich nennt deshalb das verwendete Mapping.

Der Körpereffekt übersteigt den Schichteffekt um eine Größenordnung, der Taschenkörper ist in jeder Stufe nachgiebiger als der Ring. \figref{fig:koerpereffekt_16nm_qs} zeigt den Lastschritt $16$\,Nm bei QS. Die Drehung nimmt dort bei W1 und W2 um +$14{,}6$\,\% bis +$23{,}3$\,\% zu, bei W3 und W4 um +$11{,}2$\,\% bis +$18{,}3$\,\%, an C jeweils am wenigsten. Die Druckspannung $|\sigma_{\mathrm{3}}|$ wächst in allen Stufen an A, B, D und E um +$7{,}8$\,\% bis +$18{,}9$\,\%, an C nur um $-$$2{,}7$\,\% bis +$6{,}7$\,\%. Bei der Zugspannung trennen sich die Stufen. W1 liegt an allen Punkten über dem Ring, an C um +$4{,}0$\,\% und an E um +$16{,}9$\,\%; W2 liegt zwischen +$2{,}2$\,\% und +$8{,}5$\,\%. Bei W3 liegt $\sigma_{\mathrm{1}}$ an A, C und D um $12{,}6$\,\% bis $13{,}5$\,\% unter dem Ring, an B gleich und an E um +$4{,}6$\,\% darüber. Bei W4 beträgt die Abnahme an C $8{,}6$\,\%, sonst bleibt der Unterschied innerhalb von $3{,}2$\,\%. Der anisotrope Zahn biegt sich weniger, und der Steg des Taschenkörpers trägt einen größeren Anteil. Der Kontaktdruck ändert sich bei $16$\,Nm zwischen $-$$8{,}9$\,\% (W4 an D) und +$20{,}8$\,\% (W1 an E).

\begin{figure}[htbp]
\centering
\includegraphics[width=0.8\textwidth]{kap4/koerpereffekt_16nm_qs}
\caption{Körpereffekt an den Normpunkten A bis E, Taschenkörper gegen Ringkörper gleicher Schichtung, relative Änderung je Werkstoffstufe, $16$\,Nm, quasistatisch; Größtwerte über die Zahnbreite, Marker ohne Linie, weil der Verlauf zwischen den Punkten nicht gerechnet ist.}
\label{fig:koerpereffekt_16nm_qs}
\end{figure}

\begin{table}[htbp]
\centering
\caption{Schichteffekt und Körpereffekt an den Normpunkten A bis E; Spanne der relativen Änderung je Größe und Werkstoffstufe über alle Raten, Lastschritte und Normpunkte, in \%.}
\label{tab:koerper_spannen}
\footnotesize
\setlength{\tabcolsep}{4pt}
\begin{tabular}{llllll}
\hline
Effekt & Größe & W1 & W2 & W3 & W4 \\
\hline
Schichteffekt & Ritzeldrehung & $0{,}0$ & $0{,}0$ & $-$$0{,}2$ bis $-$$0{,}1$ & $-$$0{,}1$ \\
 & Kopfverschiebung & $0{,}0$ & $0{,}0$ & $-$$0{,}2$ bis $-$$0{,}1$ & $-$$0{,}2$ bis $0{,}0$ \\
 & Kontaktdruck & $-$$0{,}1$ bis $0{,}0$ & $-$$0{,}1$ bis $0{,}0$ & $-$$1{,}8$ bis +$7{,}1$ & $-$$5{,}8$ bis +$6{,}0$ \\
 & $\sigma_{\mathrm{1}}$ & $0{,}0$ bis +$0{,}1$ & $-$$0{,}1$ bis $0{,}0$ & $-$$0{,}2$ bis +$0{,}4$ & $0{,}0$ bis +$0{,}8$ \\
 & $|\sigma_{\mathrm{3}}|$ & $0{,}0$ & $0{,}0$ & $-$$1{,}2$ bis +$1{,}1$ & $-$$1{,}3$ bis +$3{,}6$ \\
\hline
Körpereffekt & Ritzeldrehung & +$12{,}5$ bis +$23{,}3$ & +$11{,}0$ bis +$24{,}0$ & +$9{,}5$ bis +$17{,}9$ & +$8{,}3$ bis +$18{,}8$ \\
 & Kopfverschiebung & +$12{,}8$ bis +$28{,}7$ & +$11{,}9$ bis +$26{,}5$ & +$9{,}7$ bis +$18{,}5$ & +$8{,}6$ bis +$18{,}2$ \\
 & Kontaktdruck & +$1{,}3$ bis +$24{,}8$ & $-$$0{,}8$ bis +$23{,}8$ & $-$$8{,}2$ bis +$21{,}0$ & $-$$8{,}9$ bis +$15{,}0$ \\
 & $\sigma_{\mathrm{1}}$ & +$2{,}9$ bis +$16{,}9$ & +$1{,}1$ bis +$12{,}1$ & $-$$15{,}9$ bis +$4{,}8$ & $-$$14{,}1$ bis +$1{,}5$ \\
 & $|\sigma_{\mathrm{3}}|$ & +$5{,}7$ bis +$20{,}8$ & +$2{,}0$ bis +$16{,}3$ & $-$$2{,}7$ bis +$19{,}8$ & +$3{,}0$ bis +$19{,}7$ \\
\hline
\end{tabular}
\end{table}

Der Körper verändert die Verformung damit in derselben Größenordnung wie der Wechsel der Werkstoffstufe. Der Taschenkörper nach Zeichnung gilt im Weiteren als Modell des realen Teils, der Ringkörper bleibt die Bezugsgeometrie der Netzstudien aus \secref{sec:netzkonvergenz} und \secref{sec:kopfkante}. Alle weiteren Rechnungen dieser Arbeit, die Fußrundungsvarianten in \secref{sec:fussgeometrie} und die Vollserie für die Bemessung in \secref{sec:stommel_ergebnis}, laufen deshalb auf dem Taschenkörper; die anisotropen Stufen verwenden dort das Mapping auf das jeweilige Netz.

% Herkunft:
% Absatz 1 (Einleitung): series_..._qs_gerechnet/analyse/analyse.md Kopf und Abschnitt 1 (99 Stellungen, 395 von 396); fe_simulation_notes.yaml outcome von pilot_..._dy1dy2_AtoE_gerechnet (E 2543 / 3597 MPa; Raten 3 und 60 1/s aus der purpose-Zeile des Piloten); Größtwerte über die Breite: Kopf der Analyse. Stufenbezeichnungen nach Vorgabe; Zitat Sto18 für die Abstufung, Seite in 3.6 prüfen.
% Absatz 2 (Anisotropie linear): series analyse.md Abschnitt 3, Tabelle 16 Nm, Zeilen "Anisotropie (linear)" (Maximum und A bis E).
% Absatz 3 (Plastizität, W4 gegen W2): series analyse.md Abschnitt 3, Tabellen 16 Nm (Plastizität isotrop, Plastizität anisotrop, Anisotropie plastisch) und 8 Nm (sigma1 −31,9 %); Abschnitt 2, Tabelle 16 Nm (86,03 / 52,89 MPa); Fließkurve 4,3 MPa und "nahezu multiplikativ" aus ADR-116 Punkt 47(b).
% Tabelle tab:stufen_16nm: series analyse.md Abschnitt 2 (Größtwerte W1_QS, 16 Nm) und Abschnitt 3 (Spalte Maximum, 16 Nm).
% Absatz 4 (Dehnraten): pilot_..._dy1dy2_AtoE_gerechnet/analyse/analyse.md Abschnitt 8, Tabellen 16 Nm (alle Zahlen), 12 Nm (+40,2 % Kontaktdruck W3 an C, −22,2 % Drehung W4 an C), 8 Nm (Spannen über alle Lastschritte); verlängerter Eingriff aus series analyse.md Abschnitt 4 (0,47 bis 0,60 / 0,35 bis 0,48 Teilungen); Deutung der Lastaufteilung aus dem Tagebuch-outcome des Piloten ("der steifere Werkstoff teilt die Last weniger"), dort auf A bezogen, hier nach den Tabellen auf C.
% Absatz 5 (Definition Schicht- und Körpereffekt): vergleich_koerper_AtoE/analyse.md Einleitung; Schichtmaße aus known_limits.md FE-23 (0,75 gegen 0,667 / 0,786 mm) und fe_simulation_notes.yaml purpose von pilot_..._pocketlevels_..._w1w2 (6 Schichten über |z| ≤ 2,0, 7 je Taschentiefe); 60 Rechnungen je Körper, alle vollständig: outcome von pilot_..._pocket_..._gerechnet und pilot_..._pocketlevels_..._gerechnet, ADR-116 Punkt 48.
% Absatz 6 (Schichteffekt): vergleich_koerper_AtoE/analyse.md "Zusammenfassung der Effekte"; 21 % der Zahnelemente um mehr als 10° aus ADR-116 Punkt 48(a); Folgerung für die Modelldefinition aus known_limits.md FE-23.
% Absatz 7 (Körpereffekt): vergleich_koerper_AtoE/analyse.md "Zusammenfassung der Effekte" (Spannen) und Tabelle 16 Nm, Zeilen W1_QS bis W4_QS (Punktwerte); "der anisotrope Zahn biegt sich weniger, der Steg trägt mehr" aus ADR-116 Punkt 48(b).
% Bild fig:koerpereffekt_16nm_qs: bilder.yaml Eintrag koerpereffekt_16nm_qs; Unterschrift nach vergleich_koerper_AtoE/bildunterschriften.md (koerpereffekt_LOAD_16NM_QS), ohne Doppelpunkt umgestellt.
% Tabelle tab:koerper_spannen: vergleich_koerper_AtoE/analyse.md "Zusammenfassung der Effekte", alle 20 Zeilen.
% Absatz 8 (Folgerung): ADR-116 Punkt 48(b) (Körper wirkt wie der Werkstoffwechsel; Taschenkörper als Modell des realen Teils, Ring als Bezug der Netzstudien) und Punkt 53 (Entscheidung des Nutzers: jede weitere Rechnung auf dem Taschenkörper, Varianten und Vollserie, Mapping je Variantennetz).
